\documentclass[]{article}
\usepackage{lmodern}
\usepackage{amssymb,amsmath}
\usepackage{ifxetex,ifluatex}
\usepackage{fixltx2e} % provides \textsubscript
\ifnum 0\ifxetex 1\fi\ifluatex 1\fi=0 % if pdftex
  \usepackage[T1]{fontenc}
  \usepackage[utf8]{inputenc}
\else % if luatex or xelatex
  \ifxetex
    \usepackage{mathspec}
  \else
    \usepackage{fontspec}
  \fi
  \defaultfontfeatures{Ligatures=TeX,Scale=MatchLowercase}
\fi
% use upquote if available, for straight quotes in verbatim environments
\IfFileExists{upquote.sty}{\usepackage{upquote}}{}
% use microtype if available
\IfFileExists{microtype.sty}{%
\usepackage{microtype}
\UseMicrotypeSet[protrusion]{basicmath} % disable protrusion for tt fonts
}{}
\usepackage[unicode=true]{hyperref}
\hypersetup{
            pdfborder={0 0 0},
            breaklinks=true}
\urlstyle{same}  % don't use monospace font for urls
\usepackage{longtable,booktabs}
\IfFileExists{parskip.sty}{%
\usepackage{parskip}
}{% else
\setlength{\parindent}{0pt}
\setlength{\parskip}{6pt plus 2pt minus 1pt}
}
\setlength{\emergencystretch}{3em}  % prevent overfull lines
\providecommand{\tightlist}{%
  \setlength{\itemsep}{0pt}\setlength{\parskip}{0pt}}
\setcounter{secnumdepth}{0}
% Redefines (sub)paragraphs to behave more like sections
\ifx\paragraph\undefined\else
\let\oldparagraph\paragraph
\renewcommand{\paragraph}[1]{\oldparagraph{#1}\mbox{}}
\fi
\ifx\subparagraph\undefined\else
\let\oldsubparagraph\subparagraph
\renewcommand{\subparagraph}[1]{\oldsubparagraph{#1}\mbox{}}
\fi

\date{}

\begin{document}

Tutorials on computer simulations for X-ray optics

Part II: TUTORIALS on ray tracing with ShadowOui

Manuel Sanchez del Rio, ESRF, Grenoble (France)

\href{mailto:srio@esrf.eu}{\nolinkurl{srio@esrf.eu}}

(last updated 2016-01-13)

These tutorials consist on a set of examples of different practical
cases to be run using ShadowOui, the Oasys User Interface for SHADOW.

\begin{itemize}
\tightlist
\item
  11 - Geometrical source. Learning reference frames
\item
  12 - Synchrotron sources: Bending magnets
\item
  13 - Insertion devices
\item
  14 - Beam propagation (phase space (z,z') ellipses)
\item
  15 - Focusing with grazing incidence mirrors: effect of aberrations
\item
  16 - Kirkpatrick-Baez system
\item
  17 - Double crystal monochromator
\item
  18 - Sagittal focusing - python script
\item
  19 - Simulation of a complete beamline
\end{itemize}

\begin{itemize}
\tightlist
\item
  20 - Slope errors
\item
  21 - Thermal bump
\item
  22 - Curved crystal monochromators: Rowland and off-Rowland
  configurations
\item
  23 - Crystals in Laue geometry
\item
  24 - Transfocators
\item
  25 - Fresnel propagator
\item
  26 - Two slits experimnt - python scripts
\item
  27 - More examples
\end{itemize}

Appendix -- The very basics of SHADOW

\begin{itemize}
\tightlist
\item
  1 - SHADOW Introduction
\item
  2 - SHADOW files
\item
  3 - SHADOW frame
\item
  4 - Effect of optical element orientation in SHADOW frame
\item
  5 - Script programming (python): Survival guide
\item
  6 - Resources
\end{itemize}

\section{Learning reference frames in SHADOW using a geometrical
source.}\label{learning-reference-frames-in-shadow-using-a-geometrical-source.}

You will:

\begin{itemize}
\tightlist
\item
  learn to define geometrical sources
\item
  understand the use of a python script for modifying an existing source
\item
  understand reference frames
\end{itemize}

\begin{enumerate}
\def\labelenumi{\roman{enumi})}
\tightlist
\item
  Create a collimated (i.e., zero divergence) geometric source with
  elliptical shape with vertical semi axis twice the horizontal semi
  axis (e.g., 0.2 cm and 0.1 cm in Z and X, respectively). Visualize it
\item
  Apply a python script that received the source and keeps only the rays
  with positive values of X and Z (i.e., sets the flag as ``lost'' for
  rays with negative values of X and Z), and resend the beam. Visualize
  the new result after this modification of the beam.
\item
  Create a mirror optical element, with incident angle 45 deg, and
  p=q=1m. Trace the system in two cases, with Mirror orientation angle 0
  and 90 degrees. Verify the results with the pictures shown before.
\end{enumerate}

Hints: you may load the workspace ex11\_referenceframe.ows.

Answer

Pay attention to make plots using ``Rays: Good Only'' and same aspect
ratio (click the blue octagon in )

Source (x,z) plane before and after applying the script that select only
``positive'' rays

a) Mirror orientation angle = 0

Left: footprint (Y,X) plane. Right: image (X,Z) plane

b) Mirror orientation angle = 90 deg

Left: footprint (Y,X) plane. Right: image (X,Z) plane:

\section{Synchrotron sources: Bending
magnets.}\label{synchrotron-sources-bending-magnets.}

You will learn to

\begin{itemize}
\tightlist
\item
  simulate bending magnets
\end{itemize}

i) Simulate the source for the ESRF bending magnet (full emission) at a
fixed energy (e.g.,. 8 keV). Use one mrad of horizontal divergence.
Visualize the cross section (x,z), the divergence space (x',z'), the
phase spaces (x,x') and (z,z'). Visualize the top view (y,x). Make
histograms of intensity (total, -polarized and -polarized) as a
function of the vertical divergence. Plot also the degree of circular
polarization (S3 component of the Stokes vector).

ii) Change the source energy to 18 keV and compare the plot f intensity
versus vertical divergence with the result at 8 keV. Verify that the
radiation is more collimated. Simulate the same source but on a limited
vertical divergence (e.g., +/- 50 rad).

Hints: you may load the workspace ex12\_bendingmagnet.ows, where this
system is defined, for full vertical emission.

Notes:

By default one can use Calculation Mode = Precomputed, and select an
``infinite'' \emph{Max vertical half-divergence} (we set to 1 rad which
contains everything that is radiated). Shadow will calculate rays
following the full emission. In case that one wants to work very far
from the critical energy (like for infrared beamlines), one should use
``Exact Calculation''. In this case, the \emph{Max vertical
half-divergence }should be larger than the natural full divergence but
not much larger. For example, in this case one can set 0.001: 1 mrad is
still much larger than the full vertical emission (\textasciitilde{}200
mrad).

Answer

i)

Open the ``Bending Magnet'' widget, click ``Run Shadow/Source'' and
display in the right panel using ``Detailed Plot'' the cross section
(x,z), the divergence space (x',z'), and phase spaces (x,x'), (z,z'):

Another way is to plug a ``Plot XY'' widget and select the wanted
coordinates for the plot. This can also be used to display the ``top
view'' (y,x).

Plug the ``Histogram'' widget to display histograms of intensity
weighted by 23:Total Intensity, 24:-polarized, 24: -polarized,
S3-Stokes or circular polarization). Although it is not possible to
combine histograms in a single plot, it is always possible to prepare a
simple script that performs customized plots, like the one included in
the workspace that produces:

\section{Insertion devices}\label{insertion-devices}

You will learn to

\begin{itemize}
\tightlist
\item
  Simulate wigglers and undulators
\end{itemize}

a) Simulate the old wiggler for the ESRF ID17 (medical beamline) the
energy interval 10000±10 eV . Calculate the total horizontal divergence
(width of the x' histogram) and visualize a top view of the emission
(y,x) with finite emittance and source size, and without emittances
(i.e., setting emittances and sigma's to zero).

b) Simulate the ESRF U46 undulator with gap tunned to have its third
harmonic at 7833.5 eV. Use Gaussian approximation and understand the
parameters.

Hint: you may use the workspace file ex13\_insertiondevices.ows

Answer

a)

\begin{longtable}[]{@{}@{}}
\toprule
\tabularnewline
\tabularnewline
\bottomrule
\end{longtable}

Plot of X versus Y for the wiggler with (left) and without (right)
emittances

b)

Use the ``Undulator Gaussian'' widget and input the corresponding
parameters for the \emph{electron beam} and \emph{photon energy}. This
application will create a source using the Shadow Geometrical source
with divergences and sizes corresponding to the \emph{photon beam}.
These values or the photon beam (S,S') comes from a convolution (sum in
quadrature) of the values for the electron beam (s\textsubscript{e
}s\textsubscript{e}' ) with the values corresponding to the photon
emission of a single electron beam ( s\textsubscript{g},
s\textsubscript{g}') and wavelength l=1.58274326 A (7833.5 eV), and
undulator length L=1.65 m .

The formulas used are (Onuki \& Elleaume: Undulators, Wigglers and their
applications, CRC press, 2002):

In our case:

\begin{longtable}[]{@{}@{}}
\toprule
\tabularnewline
\tabularnewline
\tabularnewline
\bottomrule
\end{longtable}

These values can be checked in the ``Source Info'' available using the
``Info'' widget.

Plot of the divergence space (X',Z') for the undulator.

\section{Beam propagation (phase space (z,z')
ellipses)}\label{beam-propagation-phase-space-zz-ellipses}

You will learn to

\begin{itemize}
\tightlist
\item
  Define screens and slits associated to optical elements
\item
  Learn about the phase space changes when the beam propagates
\item
  ``Optimize the source'' in the sense shadow will only store rays that
  enter in a defined aperture
\end{itemize}

a) Using the created bending magnet source (example 12) add several
screens at 0 (source position), 50, 75 and 100 cm from the source. See
the tilt of the (z,z') diagram.

b) Define an aperture (20 m×20 m) in the first screen and see its
effect in screens \#2 and \#3. Use

Hints: You may use the workspace files ex14a\_beampropagation.ows. It
also, contains the system with slit, but using an optimised source in
order to avoid losing most of the rays at the slit.

Answer

a)

\protect\hypertarget{anchor-2}{}{}0cm 50cm

75cm 100cm

Note that instead of defining the different screens, one can directly
plug a ``Plot XY'' to the source and widget and use ``Position of the
image'' = Retraced and play with the distances.

b)

75cm 100cm

\section{Focusing with grazing incidence mirrors: effect of
aberrations.}\label{focusing-with-grazing-incidence-mirrors-effect-of-aberrations.}

You will learn to

\begin{itemize}
\tightlist
\item
  Use different mirror shapes in SHADOW
\item
  Experience with the automatic calculation of the mirror parameters.
\item
  Include mirror reflectivity by using the ``prerefl'' preprocessor
\item
  Visualize results using contour curves.
\end{itemize}

Create a geometrical source with Gaussian shape (\textsubscript{x}=57
m, \textsubscript{z}=10.4 m) and Gaussian divergence
(\textsubscript{x'}=88.5 rad and \textsubscript{z'}=7.2 rad) to
simulate the emission of an ESRF 1.65 m undulator at 10 keV in a Low
beta section.

a) Study the case of different mirror shapes (spherical, toroidal and
ellipsoidal) for focusing the source with distances (p,q)=(30m,10m)
(magnification 1/3) and (30m,1m) (magnification 1/30). Set the grazing
angle to 2 mrad. Study the effect of the spherical aberrations and its
influence depending on the magnification factor.

Study the dependence on mirror dimensions and incident angle.

b) Enter the effect of mirror reflectivity. Consider a Rh (=12.4
g/cm\textsuperscript{3}) coating and a source with energy distribution
in 5-45 keV (box-distribution). Visualize the results. Plot also the
intensity versus energy.

Hints. For the different mirrors, Shadow can calculate the surface
parameters (curvature radii, ellipse axes, etc). by selecting the
parameters in Basic Settings → Surface Shape → Type =
``internal/calculated'' . The resulting mirror parameters can be seen
using the ``MirInfo'' tab from the ``Info'' widget. For including
reflectivity, the preprocessor PreRefl must be used. The workspace file
ex15a\_aberrations.ows contains this exercise.

Answer

a) The aberrations effect increases if

\begin{itemize}
\item
  \begin{enumerate}
  \def\labelenumi{\arabic{enumi}.}
  \item
    \begin{enumerate}
    \def\labelenumii{\roman{enumii})}
    \tightlist
    \item
      one goes to more grazing angle
    \item
      one reduces the magnification factor (i.e., mode demagnification
      of the source)
    \item
      one uses larger mirrors. Small mirrors reduce aberration because
      cut rays which arrive far from the mirror center. Obviously, this
      effect reduces also the intensity. For analysing that, use the
      Basic Settings-\textgreater{}Dimensions-\textgreater{}Limits Check
      entry.
    \end{enumerate}
  \end{enumerate}
\end{itemize}

M=1/3 Toroidal: 43 mm x 26 mm FWHM, Ellipsoidal 40 mm x 9 mm

M=1/30 Toroidal: ? mm x 39 mm FWHM, Ellipsoidal 2.4 mm x 1.6 mm. Note
that a graph type ``preview'' has bee used to better observe the
aberration tails produced by the toroidal mirror.

b) I(E) plot.

\section{Kirkpatrick-Baez system}\label{kirkpatrick-baez-system}

You will learn to

\begin{itemize}
\tightlist
\item
  define an optical system with two mirrors either of circular or
  elliptical section
\item
  tell SHADOW to calculate automatically the mirror parameters in the
  case that focal planes are not coincident with continuations planes
\item
  lear the impact of the ``mirror orientation angle''
\item
  define mirror dimensions
\end{itemize}

Study the case of the previous exercise (M=1/30) with a Kirkpatrick-Baez
system with cylindrical (and later with elliptical) mirrors of length=40
cm and width=4 cm. Distance source-M1=29.5m; Distance M1-M2=1m; Distance
M2-spot=9.5m. Use 25000 rays. Do not include the mirror reflectivity.

Hint: you may use the workspace file ex16a\_kb.ows.

Answer

First branch corresponds to the KB with circular (cylindrical) mirrors.
Note that the image plane displays (X vs Z) in the ``Plot XY'' widget,
corresponding to the horizontal (Z) and vertical (X) directions. Note
that these directions are swapped respect to the source because the
mirror orientation angle for the second mirror is 90 degrees. The mirror
parameters are input externally R(M1)=739455.7 cm and R(M2)=691710.4 cm.
The Resulting spot is 36 mm x 8.3 mm. The second branch corresponds to
the KB with elliptical mirrors, where the ellipse parameters are
calculated internally. (check them using MirInfo) The Resulting spot is
\textasciitilde{}31 mm x 7 mm. A third branch uses the widget for the
compound element ``Kirkpatrick-Baez Mirror'' which reproduces the
elliptical KB setup in a simpler way. This branch also implements a
trick useful in some cases: if one wants to reverse the axes to come
back to X in horizontal and Z in vertical is possible to use an ``Empty
Element'' with incident angle zero, output angle 180 deg and mirror
orientation angle 90 deg.

\section{Double crystal monochromator.
}\label{double-crystal-monochromator.}

You will learn to

\begin{itemize}
\tightlist
\item
  create crystal reflectivity data using the ``bragg'' preprocessor
\item
  use the ``autotunning'' facility to align the crystal
\item
  calculate the energy resolution for a crystal and a combination of
  systems
\item
  optimise the source bandwidth
\item
  play with the mirror orientation angle. Relate its values to the
  crystal dispersion ( (+,-) and (+,+) crystal combination)
\end{itemize}

Create a bending magnet source (starting from exercise 12) at 8000±25 eV
with 3 mrad horizontal divergence. Verify its energy dependence and
horizontal and vertical divergence values.

a) Implement a flat Si 111 crystal at 30 m from the source. Verify the
energy dependence and calculate resolution. Redefine the source energy
bandwidth to optimize the calculation in order to obtain the energy
dependence with the highest signal.

b) Add a second crystal 10 cm downstream from the first one in (+,-) and
(+,+) configurations (play with the mirror orientation angle). Explain
the obtained differences in energy resolution.

Hint: you may use the workspace file ex17\_crystalmono.ows containing
three branches: i) (+,-) ii) (+,+), and iii) the double crystal
monochromator compound element (same as (+,-). Use the Bragg
preprocessor to create the reflectivity data for a Si 111 crystal. You
can create the output file for a large range of energy (e.g., from 5000
to 15000 eV). Pay attention to the file name when you run bragg, because
it also appears in the o.e. crystal menus. In this case, the file is
called si5\_15.111.

Answer

a) The optimized energy range selected is 8000±10 eV. The histogram of
the energy (including reflectivity) after the first crystal has
FWHM\textasciitilde{}5.6eV:

b) Left: resolution function for (+,-) (non-dispersive configuration).
Here the mirror orientation angles are 0 and 180 deg for the first and
second crystals, respectively. Note that the width (5.2 eV) is very
similar to a single crystal in a), but with a slightly lower intensity
due to the absorption of the second crystal. Right: resolution function
for (+,+) (dispersive configuration). The mirror orientation angle is 0
for both crystals. Here one can see the effect of the dispersive setup,
the final energy resolution depends only on the crystal Darwin width and
does not depend on the beam divergence.

\section{Sagittal focusing - python
script}\label{sagittal-focusing---python-script}

You will learn to

\begin{itemize}
\tightlist
\item
  define a cylindrical mirror for sagittal focusing
\item
  define ``externally'' the optical element radius of curvature
\item
  optimize the focal spot
\end{itemize}

a) Using the (+,-) system defined in the last exercise at photon energy
20 keV, bend sagittally the second crystal to focus in the horizontal
plane at the sample position, placed 1000 cm downstream from the
monochromator (monochromator at 3000 cm from the source). Calculate
horizontal spot size.

b) Study the effect of the ratio between the distances mono-sample and
source-mono in the transmitted intensity. Study the case of M=1/30. See
the effects in energy resolution and system transmitivity. Explain these
differences. Verify that ratio 1/3 is the optimum.

Hints: You may use the workspace files ex18a\_sagittalfocusing.ows also
contains a python script that scans the magnification. One can see that
the intensity peaks at about 1/3. A script is included to calculate the
curvature radius: Rs (20 keV, M=1/3)=148.3 cm.

Answer

a) Focusing system: Spot width H=334 mm,
I/I\textsubscript{0}=1090/25000, E=18 eV

Non-focusing system: Spot width H=189mm,
I/I\textsubscript{0}=1094/25000, E=18 eV

Best focus very close to the focal position. The spot size does not
change appreciably.

b) p=30m, q=1m, R\textsubscript{s}=47.8 cm, Spot width=0.012 cm,
I/I\textsubscript{0}=168/25000, E=5.8 eV

The study of the variation of the intensity as a function of the
magnifications needs to run shadow for many points of M. It can be done
with a macro. The result should show an optimum magnification of M=1/3
for large divergence values. The python script will show this effect (5
mrad).

Intensity (in arbitrary units) versus magnification factor M for a point
and monochromatic (E= 20 keV) source placed at 30 m from the sagittaly
bent crystals. The beam divergences is 5 mrad (can be changed in the
macro). We clearly observe the maximum of the transmission at M=0.33, as
predicted by the theory.

\section{Simulation of a complete
beamline.}\label{simulation-of-a-complete-beamline.}

You will learn to:

\begin{itemize}
\tightlist
\item
  Combine several optical elements
\item
  Obtain final results for a beamline in terms of flux, resolution and
  spot size. 
\end{itemize}

Define the following elements in SHADOW:

Geometrical Gaussian source at 10000±10 keV (box distribution) (like in
exercise 15 b, but changing the energy interval)

M1: Cylindrically collimating mirror in the vertical plane at 25 m.
Grazing angle 0.12 degrees. Rh coating (density=12.4 g/cc). Infinite
dimensions.

MONO: Double crystal monochromator, Si 111, with second crystal
sagittally bent (focusing the source into the sample position in the
horizontal plane), at 30 m from the source (Rs=296.6 cm)

M2: Re-focusing mirror at 35m from the source, focusing at the sample
position. Same angle as M1

Sample at 40 m

Calculate:

\begin{enumerate}
\def\labelenumi{\roman{enumi})}
\tightlist
\item
  Beam geometry at the sample position
\item
  Energy resolution
\item
  Transmitivity of the whole beamline. Number of photons at the sample
  position supposing that at the source we have, at 10 keV, a flux of 5
  10\textsuperscript{13} ph/sec/0.1\%bw
\end{enumerate}

b) How are these results modified using a focusing first mirror and a
flat second mirror?

Hint: you may use the workspace file ex19\_beamline.ows

Answer

\textsubscript{source}=4 eV (optimized source bandwidth); =1.3eV;
I/I\textsubscript{0}=1406/5000

Transmitivity in one eV=T=(I/
/(I\textsubscript{0}/\textsubscript{source})=(1406/1.3)/(5000/4)

Number of photons at the source in one eV bandwidth =N= 5
10\textsuperscript{12}

Total number of photons = N×T×E

Intensity distribution in the (X,Z) plane at the image position

Energy distribution:

\section{Slope errors.}\label{slope-errors.}

You will learn to:

\begin{itemize}
\tightlist
\item
  Use Waviness preprocessor to create a file sampling slope errors
\item
  Use presurface to inject it in SHADOW
\item
  See the important effect of slope errors in the focal size 
\end{itemize}

a) Load the Kirkpatrick-Baez system of exercise 16. Set mirror surface
to be elliptical. Check that mirror dimensions are 40×4
cm\textsuperscript{2}. Calculate spot sizes without slope errors.

b) use the preprocessor widget Waviness to create maps of slope errors.
For adjusting the slope errors to the wanted values, modify the value of
the initial Y slope error in order to get a value close to the desired
tangential slope error of 0.5 arcsec rms. Then modify the number of
points in X in order to adjust the sagittal slope error to 1 arcsec rms.
In the oe ``Advanced setting'' tab, select ``Modified surface'' subtab
and check the file name containing the errors. This is automatically
populated if one connects the Waviness with the Mirror widgets. One can
also have a quick preview from there/

Hints: use the workspace ex20\_slopeerrors.ows

Answer

Top: No slope errors: 7.5 (V) × 40 (H) m\textsuperscript{2}. Bottom:
with Slope errors: 111 (H) × 89 (V) m\textsuperscript{2}

\section{Thermal bump.}\label{thermal-bump.}

You will learn to:

\begin{itemize}
\tightlist
\item
  use a python script to create a file sampling a thermal bump
\item
  see the effect of the bump in energy resolution. 
\end{itemize}

Load the ex21\_thermalbump.ows workspace. Run the script to create a
Gaussian bump bump.dat Run the system (a single Si111 crystal) without
and with thermal bump. See the changes in the energy resolution.

Answer

No bump: 1.4 eV, H: 8.5 mm, V: 0.59 m

With bump: 9.6 eV, H: 10.1 mm, V: 4.4 m

\section{Curved crystal monochromators: Rowland and off-Rowland
configurations}\label{curved-crystal-monochromators-rowland-and-off-rowland-configurations}

You will learn to:

\begin{itemize}
\tightlist
\item
  understand the effect of crystal radius in energy resolution and
  focusing conditions
\item
  calculate the focusing conditions in and out Rowland configuration 
\item
  understand the importance of using contour curves with PlotXY
\item
  simulate an asymmetric crystal
\end{itemize}

a) Using the same Gaussian source as in exercise 21, verify the focusing
conditions for a symmetrical Si111 Bragg crystal at 10 keV, with p=30m.
Calculate E. Calculate E for R\textsubscript{t}=5000 cm and
R\textsubscript{t}=2500 cm. Explain the differences.

b) Calculate the Rowland conditions for 10 keV, Si111, p=30m and
asymmetry angle =5º. Calculate energy resolution and spot size.

Hint: use ex22\_rowland.ows

Answer

a) For the focusing conditions (first branch), we have
R\textsubscript{t}=15171 cm (see results in Info widget) and E=1.32 eV

For R\textsubscript{t}=5000 cm (second branch) we then have E=1.38 eV

For R\textsubscript{t}=2500 cm (third branch) we then have E=3.9 eV

b) Using the script to calculate Rowland conditions, we get
R\textsubscript{t}=10623.284 cm, q=1184.8 cm; We get .E=0.72 eV; spot
size = 7.1 mm ×276 mm.

\section{Crystals in Laue geometry}\label{crystals-in-laue-geometry}

You will learn to:

\begin{itemize}
\tightlist
\item
  Set Laue crystals in SHADOW
\item
  See the transformation in the phase space
\item
  Apply a macro to copy intensity from one file to another.
\end{itemize}

Load the system from file ex23\_laue.ows, consisting in an asymmetric
Laue diamond (111) crystal (a=3.55 Å) and a symmetric Bragg germanium
(220) crystal (a=5.57 Å) in non-dispersive configuration. Show the
energy histograms after the two crystals. Study how the Laue crystal
change of phase space of the beam. Relate these changes to the Liouville
theorem.

Answer

Left: Energy resolution after the first Laue crystal (0.8 eV) and,
Right: energy resolution after the second Bragg Ge crystal (0.64 eV)

Phase space (Z,Z') after the Laue crystal:

Phase space (Z,Z') before the crystal.

Left: It can be displayed using the Plot XY widget connected to the
Source and retraced 40 m.

Right: However, it is interesting to display only the beam that will be
diffracted by the Laue crystal. For that we need to manipulate the beam
incoming to the Laue crystal, and change the ray intensities (i.e., the
electric fields) copying the values after the diffraction. This is done
with a script. It guaranteed that the area of the phase space
(corresponding to the beam intensity) is conserved as required by the
Liouville theorem.

\section{Transfocators}\label{transfocators}

You will learn to:

\begin{itemize}
\tightlist
\item
  Run a transfocator
\item
  Find the position of the best focus
\item
  Study the changes due to the source energy (chromatic aberrations)
\end{itemize}

Implement a monochromatoc Gaussian source of 48.2 (H) x 9.5 (V) m RMS
size and 100 (H) x 4.3 (V) rad RMS divergence. Create the same source
for two different energies (E1=35200 eV and E2=35700 eV).

Implement a transfocator consisting of two CRLs, 2D-focusing (in both H
and V), the first made in Be (16 lenses of radius 200 m, separated
50m), and the second in Al (21 lenses of radius 200 m, separated
50m). The focal distances are p=3150 cm and q=1000 cm. This setup is
discussen in Baltser et al.
\href{http://dx.doi.org/10.1117/12.893343}{\emph{\emph{http://dx.doi.org/10.1117/12.893343}}}

Compare the results using plotxy. Study the beam evolution close to the
focal plane using focnew and ray\_prop: you will find a small
astigmatism, i.e., the H and V foci are at slightly different position
(you will find that the astigmatism disappears when using a point
source, so it is the result of the final dimensions of the source).
Change the energy of the source from 35200 eV to 35700 eV and see the
effect in focal position and intensity.

Hint: use the system in file ex24\_transfocator.ows.

Answer

For E=35200 eV we find a focus of 34 x 10 m (FWHM) with I= 11528, best
H focus at about -40 cm (you can use focnew or ray\_prop implemented in
two scripts).

Sagittal focus at : -65.4294

Tangential focus at : -40.9559

Circle of least confusion : -54.2916

For E=35700 eV we find a focus of 34 x 6.1 m (FWHM) with I= 11734

Sagittal focus at : -32.9289

Tangential focus at : -5.03177

Circle of least confusion : -20.3752

\section{Fresnel propagator}\label{fresnel-propagator}

You will learn to:

\begin{itemize}
\tightlist
\item
  Compare the ray tracing results versus wave optics results for an
  aperture in 1D
\item
  Use a python script to compute the Fresnel wave optics propagator. 
\end{itemize}

Implement a monochromatoc (E=11000 eV) , collimated in H and with enough
divergence in V to illuminate a 100 m slit placed at 3760 cm from the
source. Compare the results by ray tracing propagation at 550 cm from
the slit with the diffraction pattern calculated by wave optics.

Hints: Use the ex25\_fresnel.ows workspace.

Answer

A monochromatic (E=11000 eV) point geometric source with negligible
horizontal divergence is created. A slit is placed and the histogram of
Z retraced at 550 cm shows a mostly planar distribution (ray tracing
results), so the projection of the screen in the image plane:

the script included in the workspace implements a Fresnel propagator
that computes the propagation of the shadow rays at the exit of the
screen, into a detector plane at 550 cm:

\section{Two slits experiment - python
scripts}\label{two-slits-experiment---python-scripts}

You will learn to:

\begin{itemize}
\item
  Reproduce experimental results of a diffraction from a double-slit
  Leitenberg \emph{et al.} Physica B 336 (2003) 63-67
  \href{http://dx.doi.org/10.1016/S0921-4526(03)00270-9}{\emph{\emph{http://dx.doi.org/10.1016/S0921-4526(03)00270-9}}}
\item
  Define a doube-slit in SHADOW using a screen/slit with external file
  definition
\item
  Define a source optimised to illuminate into a reduced acceptance.
\item
  Implement a Fresnel-Kirchhoff propagator in a python script to compute
  the diffraction patter produced by two slits
\end{itemize}

Implement a monochromatic (E=14000 eV) rectangular source of 0 (H) x 140
(V) m RMS size and 0 (H) x 2500 (V) rad RMS divergence optimised to
illuminate a square of 100 m at 3090 cm (defined in file
acceptance.dat). Trace the source into a screen containing the double
slit, defined in a file twoslits.pol. Perform Fresnel propagation to to
plane at 500 cm from the slits plane (fully coherent illumination). Do
the same for an ensemble of sources with different phases, and add the
final intensities (partial coherence).

Tips:

Use the file ex26\_two\_slits.ows

Use a source of 140 microns vertical, with enough divergence to fully

illuminate the slits plane (15.4 microns) at 30.9m, so the divergence is

(70+7.7)*1e-3/30.9 rads.

Use variance reduction at the source to improve efficiency, with
acceptance file: acceptance.dat

The two slits can be defined using external file in the Screen-Slit
widget twoslits.pol in the Screen-Slit widget

After running SHADOW, the python script performs the propagation in
vacuum and the ensemble average. Please note that there are parameters
hard-coded in this

script.

Answer

\section{More Examples}\label{more-examples}

Other workspaces with more examples are also available:

\begin{itemize}
\item
  CRL\_Snigirev\_1996.ows 
\item
  EBS\_3\_pole\_wiggler.ows 
\item
  crystal\_analyzer\_diced.ows 
\item
  crystal\_asymmetric\_backscattering.ows 
\item
  crystal\_mosaic.ows 
\item
  hybrid-mirror.ows 
\item
  hybrid-twoslits.ows 
\item
  knife\_edge\_scan.ows 
\item
  screen\_pattern.ows 
\end{itemize}

Appendix - The very basics of SHADOW

1 SHADOW introduction

SHADOW is a ray-tracing program specially optimized for the design of
the synchrotron radiation beamline optics.

SHADOW generates and traces \emph{a beam }along the \emph{optical
system}. The beam is a collection of \emph{rays }in a given point of the
beamline which are stored in a disk file. The optical system is a
collection of \emph{optical elements (o.e.) }(mirrors, multilayers,
slits, screens, etc.) placed in a sequential order.

Each ray is an array of 18 variables or \emph{column}s. Each variable of
column has an special physical meaning. The first six defines the
geometry: spatial coordinates (Col. 1,2,3 or x, y and z, respectively)
and the direction of the ray (cols. 4,5,6, or x',y' and z',
respectively). The rest of the columns defines the history of the ray
traversing the optical system (electric vector for s-polarization (cols.
7,8,9) and p-polarization (cols. 16-18), flag for lost ray (10),
wavelength (11) etc.).

The \emph{source }is the beam at the starting point. It is generated by
SHADOW by sampling the spatial, angular, energy and other qualities of
the synchrotron radiations sources (i.e., bending magnets, wigglers and
undulators) into a finite number of rays, using a Monte Carlo method. At
the source position the intensity of each ray (or better, its
probability of observation) is set to 1. This intensity will decrease
along the beamline because of the interaction of the ray with the
optical elements. The source generated by SHADOW samples linearly the
real source, which allows scaling the intensity with the number of
photons.

SHADOW traces the source sequentially thought each individual optical
element of the optical system. SHADOW solves the intercept of each ray
at a given o.e., calculates the output direction and the decrease in
intensity. This decrease is calculated for each ray using a physical
model (i.e. Fresnel equations for mirrors, Dynamical Theory of the
Diffraction for perfect crystals, etc.)

2 SHADOW files (it can be ignored for most ShadowOui users)

By default, ShadowOui runs SHADOW without writing files. It is however
possible to write these files, and to understand their meaning, in
particular when one wants to recover old SHADOW runs and export results
to other programs.

Important input files (written by SHADOW):

\begin{itemize}
\tightlist
\item
  start.xx an ASCII file with the list of variables for the source or
  optical elements (start.00 for the source, start.01 for the first o.e,
  start.02 for the second, and so on)
\item
  systemfile.dat a small file listing the start.xx files to be traced.
\end{itemize}

After running SHADOW, the binary files containing the rays at different
points are:

\begin{itemize}
\tightlist
\item
  begin.dat binary file containing the beam at the source position
\item
  mirr.xx binary file containing the beam on each o.e. (i.e. mirr.02 is
  the beam on the second o.e)
\item
  star.xx binary files with the beam at the image created by each o.e.
  The image of a given o.e. is the source for the following o.e.
\item
  screen.xxyy inary files for screens or slits. Screens allow to define
  apertures (slits or beam stoppers) and absorbers (filters). xx refers
  to the o.e. and yy refers to the screen order (i.e. screen.0204 means
  the fourth slit associated to the 2\textsuperscript{nd} o.e.)
\end{itemize}

3 SHADOW frame

The coordinate system of SHADOW is (schematic, ):

Note that:

\begin{itemize}
\tightlist
\item
  The y(2) coordinate is along the beam direction
\item
  The frame is rotated if one o.e. is rotated
\item
  The position, orientation, etc. of any o.e. is always referred to the
  previous one
\end{itemize}

4 Effect of o.e. orientation in SHADOW frame

5 Script programming (python): survival guide

Before anything, in python:

\textgreater{}\textgreater{}\textgreater{} import Shadow

There are only three main objects in Shadow python: a container for the
variables needed to create a source (the same variables as in the file
start.00), a container for the variables necessary to define an optical
element (the same as in the files start.01, start.02, etc.), and a
container for the beam (same information as in the binary files
begin.dat, star.xx and screen.xxyy)).

Initialize them as:

src = Shadow.Source()

oe1 = Shadow.OE()

oe2 = Shadow.OE()

beam = Shadow.Beam()

In the case we want to read variables from existing files do:

src.load(``start.00'')

oe1.load(``start.01'')

oe2.load(``start.02'')

For applying to the beam the source:

beam.genSource(src)

For tracing the two optical elements:

beam.traceOE(oe1)

beam.traceOE(oe2)

Access SHADOW beam

The rays are in a numpy array:

\textgreater{}\textgreater{}\textgreater{} print(beam.rays)

{[}{[} 0.00039945 0. 0.00034409 ..., 0. 0. 0. {]}

{[}-0.00109609 0. 0.00174768 ..., 0. 0. 0. {]}

{[}-0.0013806 0. -0.00160363 ..., 0. 0. 0. {]}

...,

{[}-0.0005546 0. -0.00137695 ..., 0. 0. 0. {]}

{[} 0.00045102 0. -0.00234013 ..., 0. 0. 0. {]}

{[}-0.00062769 0. 0.00102393 ..., 0. 0. 0. {]}{]}

Write binary file:

beam.write('star.02')

Visualizing results

Shadow.ShadowTools.plotxy(beam,1,3)

Shadow.ShadowTools.histo1(beam,1)

Using Scripts in ShadowOui

It is possible to connect a ``Python Script'' widget to a source or
optical element and extract and modify the information there. This is
very useful for making things that are not implemented as widgets. The
SHADOW objects are inside a large object that the ``Python Script''
widget receives as in\_object and resends as out\_object. For example:

beam = in\_object.\_beam

print((in\_object.history{[}1{]}.\_shadow\_oe\_end.\_oe.FILE\_SCR\_EXT))

6 Resources

SHADOW papers:

\begin{itemize}
\tightlist
\item
  F. Cerrina and M. Sanchez del Rio "Ray Tracing of X-Ray Optical
  Systems" Ch. 35 in Handbook of Optics (volume V, 3rd edition), edited
  by M. Bass, Mc Graw Hill, New York, 2009. ISBN: 0071633138 /
  9780071633130
  \href{http://www.mhprofessional.com/handbookofoptics/vol5.php\%20}{\emph{http://www.mhprofessional.com/handbookofoptics/vol5.php}}
\item
  M. Sanchez del Rio, N. Canestrari, F. Jiang and F. Cerrina "SHADOW3: a
  new version of the synchrotron X-ray optics modelling package" J.
  Synchrotron Rad. (2011). 18, 708-716
\end{itemize}

\href{http://dx.doi.org/10.1107/S0909049511026306\%20}{\emph{http://dx.doi.org/10.1107/S0909049511026306}}

Code repositories:

\begin{itemize}
\item
  ShadowOui-Tutorials (this document and the workspace files used here):
  \href{https://forge.epn-campus.eu/projects/shadow3/}{\emph{\emph{https://github.com/srio/ShadowOui-Tutorial/}}}
\item
  ShadowOui repository
  \href{https://github.com/srio/shadow3/}{\emph{https://github.com/lucarebuffi/ShadowOui/}}

  SHADOW3 repository
  \href{https://github.com/srio/shadow3/}{\emph{https://github.com/srio/shadow3/}}
  (check the README files in the ``doc'' folder) 
\end{itemize}

\end{document}
